\documentclass{ximera}
\title{A Ferris Wheel}

\newcommand{\pskip}{\vskip 0.1 in}

\begin{document}
\begin{abstract}
Relating small changes as you ride ferris wheel.
\end{abstract}
\maketitle




\begin{question}  \label{Qererfeft4tr4}

This question is about a ferris wheel. The idea is this. After riding a ferris wheel for some time, the wheel gets stuck. Then it turns a through a small angle and gets stuck again. The question is to relate the small change in your height to the small change in the distance you move along your circular path as the wheel turns through this small angle.

The relationship between these small changes depends upon your position on the ferris wheel. The worksheet below illustrates this. Drag the slider $m$ in Line 2 to see how the relationship between the small change $\Delta h$ in your height and the small change in the distance $\Delta s$ along your path changes as your position changes.

\begin{onlineOnly}
    \begin{center}
\desmos{vc1xpi5mbq}{450}{600}  
\end{center}
\end{onlineOnly}

\href{https://www.desmos.com/calculator/vc1xpi5mbq}{151: Ferris Wheel 57}

\begin{enumerate}
\item Summarize your observations.
\begin{freeResponse}
\end{freeResponse}


\item The first step in relating these small changes is to find a function
\[
      h = f(s) \, , s\geq 0,
\]
that expresses your height (in feet) above the ground in terms of the distance (in feet) along your path since you boarded the wheel. 

For this, we'll assume the wheel has a radius of $30$ feet and that its center is $40$ feet above the ground.




\end{enumerate}
 


The function
\[
      h = f(t)  \, , \, 0\leq t \leq 2.2 ,                     %= -2t^3+7t^2-8t+8 \, , \, 0\leq t \leq 2.2 ,
\]
expresses the height of a balloon (in hundreds of feet) in terms of the number of minutes past noon.

The graph of the function $h=f(t)$ is shown below.

\begin{onlineOnly}
    \begin{center}
\desmos{uolmkpi5yr}{450}{600}  
\end{center}
\end{onlineOnly}

\href{https://www.desmos.com/calculator/uolmkpi5yr}{151: Ferris Wheel 3}


\end{question}

\end{document}
